\tocsection{\textbf{11/9/26 Lecture 30:} Quotients, Products, and Compactness}

\textbf{Recap:} When given a sub-basis, we can make a topology using unions of finite intersections of elements of the sub-basis. When given a basis, we can do the same thing with just unions.\\

Moving onto homeomorphisms, consider the following example.

\begin{example}{\textbf{Example:} A Continuous Bijection That Isn't a Homeomorphism}
    Let $f : [0, 1) \to S^1$ by $f(x) = e^{2\pi ix}$. $f$ is continuous and bijective, but its inverse is not continuous. This is because $[0, 0.5)$ is open, but the arc $[1, -1) \subseteq S^1$ is not.
\end{example}

\subsection*{Quotient Topology}
\sectionspace

Let $(X, \cal{T})$ be a topological space and $\sim$ an equivalence relation on $X$. Now, consider
\[ X/{\sim} = \{[x] \col x \in X\} \]
We know there exists a canonical quotient map $q : X \to X/{\sim}$ by $x \mapsto [x]$. We define the topology on $X/{\sim}$ to be given by
\[ U \subseteq X/{\sim} \tt{ is open} \iff q^{-1}(U) \tt{ is open in } X \]

\begin{example}{\textbf{Example:} $\R/\Z$}
    Consider $\R$ with $x \sim y \iff x - y \in \Z$. Then, we see
    \[ \R/{\sim} = \{[x] \col x \in [0, 1)\} \]
    Another useful way to think of it is $[0, 1]$ but we identify $0$ and $1$. Note that this gives $[0, 1)$ a topology different from the subspace topology, since $[0, 0.5)$ is not open in this topology.
\end{example}

\begin{lemma}[\textbf{Proposition 30.1}]
    Let $X$ be a topological space with equivalence relation $\sim$. Let $f : X \to Y$ be continuous and constant on all equivalence classes. Then, there exists a continuous map $\tilde{f} : X/{\sim} \to Y$ such that $f = \tilde{f} \circ q$.
    \begin{center}
        \begin{tikzcd}
            X \arrow[d, "q"'] \arrow[dr, "f"] & \\
            X/{\sim} \arrow[r, "\tilde{f}"'] & Y
        \end{tikzcd}
    \end{center}
\end{lemma}
\begin{proof}{Proof:}
    Let $\tilde{f}([x]) = f(x)$, which is well defined as $f$ is constant on each equivalence class. Now, if $U$ is open in $Y$, then
    \[ q^{-1}(\tilde{f}^{-1}(U)) = f^{-1}(U) \]
    which is open as $f$ is continuous, so $\tilde{f}^{-1}(U)$ is open in $X/{\sim}$, since $q^{-1}(V)$ is open iff $V$ is open. 

    $\qed$
\end{proof}

\begin{example}{\textbf{Example:} $\R/\Z \cong S^1$}
    If $f(x) = e^{2\pi ix}$, it respects $x \sim y \iff x - y \in \Z$. Thus, we have a map $\tilde{f} : \R/\Z \to S^1$ that turns out to be a homeomorphism.
\end{example}

\subsection*{Product Topology}
\sectionspace

Let $X = \prod_{\alpha \in I} X_{\alpha}$ where each $X_{\alpha}$ is a topological space. We want to define a topology on $X$. To motivate, $\R^2 = \R \times \R$ is generated by the basis $U \times V$ ($U \subseteq \R$, $V \subseteq \R$ open).

\begin{definition}[\textbf{Definition 30.2:} Box Topology]
    The \emph{box topology} on $X$ is generated by the basis
    \[ \left\{\prod_{\alpha} U_{\alpha} \mid U_{\alpha} \subseteq X_{\alpha} \tt{ is open}\right\} \]
\end{definition}

This is not ideal, however, and we will show why.

\begin{example}{\textbf{Example:} The Issue with the Box Topology}
    Say we have $f_{\alpha} : Y \to X_{\alpha}$. We can form $f : Y \to X$ by $f(y) = \{f_{\alpha}(y)\}_{\alpha \in I}$. One desired property to want is that $f_{\alpha}$ continuous $\implies f$ continuous. It turns out the box topology doesn't give this.\\

    Let $X = \prod_{n = 1}^{\infty} \R$ with the box topology. Let $f_n(x) = x$, $f_n : \R \to \R$. Now, this gives us a map $f : \R \to X$ by
    \[ f(x) = (x, x, \ldots) \]
    It turns out this is not continuous! Consider
    \[ U = (-1, 1) \times \left(-\frac{1}{2}, \frac{1}{2}\right) \times \left(-\frac{1}{3}, \frac{1}{3}\right) \times \cdots \]
    This is open, but
    \[ f^{-1}(U) = \{0\} \]
    so $f$ is not continuous!
\end{example}

Now, given $X = \prod_{\alpha} X_{\alpha}$, we get a map $\pi_{\alpha} : X \to X_{\alpha}$ by $\pi_{\alpha}(x) = x(\alpha) := x_{\alpha}$.

\begin{definition}[\textbf{Definition 30.3:} Product Topology]
    The \emph{product topology} is the topology generated by the sub-basis
    \[ \bigcup_{\alpha \in I}\{\pi_{\alpha}^{-1}(U) \col U \subseteq X_{\alpha} \tt{ open}\} \]
    Note
    \[ \pi_{\alpha}^{-1}(U) = U \times \prod_{\alpha' \neq \alpha} X_{\alpha'} \]
    This topology can also be described as being generated by the basis
    \[ \left\{\prod_{\alpha} U_{\alpha} \col U_{\alpha} \subseteq X_{\alpha} \tt{ is open and equal to } X_{\alpha} \tt{ for all but finitely many } \alpha\right\} \]
\end{definition}

Note that the product and box topologies are the same on finite products. However, in general, the box topology is larger, or ``finer''.

\begin{definition}[\textbf{Definition 30.4:} Finer and Coarser Topologies]
    Let $\cal{T}$ and $\cal{T}'$ be topologies on the same space $X$. Suppose $\cal{T} \subseteq \cal{T}'$. We say $\cal{T}'$ is \emph{finer} than $\cal{T}$ and $\cal{T}$ is \emph{coarser} than $\cal{T}'$.
\end{definition}

\begin{lemma}[\textbf{Proposition 30.5}]
    Let $f : A \to \prod_{\alpha} X_{\alpha}$ such that $f(a) = (f_{\alpha}(a))_{\alpha}$. Then, $f$ is continuous with respect to the product topology iff $f_{\alpha}$ is for all $\alpha$.
\end{lemma}
\begin{proof}{Proof:}
    If $f$ is continuous, then as $f_{\alpha} = \pi_{\alpha} \circ f$, we see $f_{\alpha}$ is continuous.\\

    Conversely, if $f_{\alpha}$ is continuous for all $\alpha$, let
    \[ U = U_{\alpha_1} \times \cdots \times U_{\alpha_n} \times \prod_{\alpha \neq \alpha_i} X_{\alpha} \]
    be a basis element. Then,
    \[ f^{-1}(U) = \bigcap_{i = 1}^{n} f_{\alpha_i}^{-1}(U_{\alpha_i}) \]
    which is open. Thus, as the preimages of basis elements are open, $f$ is continuous.

    $\qed$
\end{proof}

\subsection*{Compactness}
\sectionspace

\begin{definition}[\textbf{Definition 30.6:} Open Covers]
    Let $(X, \cal{T})$ be a topological space. An \emph{open cover} of $X$ is a collection $\{U_{\alpha}\}$ of open sets such that $X = \bigcup_{\alpha} U_{\alpha}$. A \emph{finite subcover} is a finite subcollection $\{U_{\alpha_i}\}_{i = 1}^{n}$ such that $X = \bigcup_{i = 1}^{n} U_{\alpha_i}$.
\end{definition}

\begin{definition}[\textbf{Definition 30.7:} Compactness]
    A topological space $(X, \cal{T})$ is \emph{compact} if every open cover has a finite subcover.
\end{definition}

% NOTE (check): here the U_alpha are open sets of X (their traces U_alpha cap A are the open sets of A), and it's "iff": A is compact in the subspace topology exactly when every such cover of A has a finite subcover.
\textbf{Note:} If $A \subseteq X$, $A$ is compact in the subspace topology iff
\[ A \subseteq \bigcup_{\alpha} U_{\alpha} \implies A \subseteq \bigcup_{i = 1}^{n} U_{\alpha_i} \]

\begin{lemma}[\textbf{Proposition 30.8}]
    A closed subset of a compact set is compact.
\end{lemma}
\begin{proof}{Proof:}
    Let $A \subseteq X$ be closed. Then, if $\{U_{\alpha}\}$ is an open cover of $A$, then $\{U_{\alpha}\} \cup \{A^{c}\}$ is an open cover of $X$. Now, we get a finite subcover
    \[ X = \bigcup_{i = 1}^{n} V_{\alpha_i} \implies A \subseteq \bigcup_{i = 1}^{n} V_{\alpha_i} \]
    % FIXED: page has the V's drawn from {U_alpha} cup {A}; changed A to A^c (that's the set added to make the cover of X, and the one removed next).
    where $V_{\alpha_i}$ is drawn from $\{U_{\alpha}\} \cup \{A^{c}\}$. Removing $A^{c}$ from the collection, we still cover $A$ with a finite subcollection of $\{U_{\alpha}\}$.

    $\qed$
\end{proof}
